\documentclass[11pt]{article}

\usepackage{amsmath}
\usepackage{amssymb}
\usepackage{hyperref}
\usepackage{tikz}

\title{
    Intelligent System 2\\
    Tutorial week 11: Abstract Argumentation\\
    \textbf{Solutions}
}

\author{
    Matteo Magnini\\
    \href{mailto:matteo.magnini@uni.lu}{matteo.magnini@uni.lu}
}

\date{
    6 May, 2026
}

\begin{document}

\maketitle

\section{Solutions}

\subsection{Exercise 1}

\textbf{Recall:} A set is conflict-free iff no arguments in it attack each other.

Attacks:
\[
a \leftrightarrow b,\quad b \to c,\quad c \leftrightarrow d
\]

\textbf{Conflict-free sets:}
\[
\emptyset,\ \{a\},\ \{b\},\ \{c\},\ \{d\},\ \{a,c\},\ \{a,d\},\ \{b,d\}
\]

\textbf{Maximal (naive candidates):}
\[
\{a,c\},\ \{a,d\},\ \{b,d\}
\]

These are maximal because adding any other argument introduces a conflict.

\subsection{Exercise 2}

\textbf{Naive extensions = maximal conflict-free sets:}
\[
\{a,c\},\ \{a,d\},\ \{b,d\}
\]

\textbf{Example non-naive:}
\[
\{a\}
\]
It is conflict-free but not maximal (can be extended to $\{a,c\}$).

\subsection{Exercise 3}

\textbf{Recall:} $E$ defends $a$ if for every attacker $b \to a$, there exists $c \in E$ such that $c \to b$.

Attacks:
\[
b \to a,\quad c \to b,\quad d \to c
\]

\begin{itemize}
    \item $\{a\}$:
    $a$ is attacked by $b$, but $\{a\}$ does not attack $b$ $\Rightarrow$ not defended.

    \item $\{c\}$:
    $c$ is attacked by $d$, but $\{c\}$ does not attack $d$ $\Rightarrow$ not defended.
\end{itemize}

\textbf{Admissible sets:}

\[
\emptyset,\ \{d\}
\]

Explanation:
\begin{itemize}
    \item $\emptyset$: trivially admissible
    \item $\{d\}$: conflict-free and has no attackers
\end{itemize}

\subsection{Exercise 4}

\textbf{Recall:} Complete = admissible + contains all defended arguments.

Admissible sets:
\[
\emptyset,\ \{d\}
\]

\textbf{Check completeness:}

\begin{itemize}
    \item $\emptyset$:
    does not defend any argument $\Rightarrow$ complete

    \item $\{d\}$:
    $d$ has no attackers, so it is trivially defended.

    However, $\{d\}$ also defends $b$:
    \[
    c \to b \quad \text{and} \quad d \to c
    \]
    hence $d$ attacks all attackers of $b$.

    Therefore, $\{d\}$ defends $b$ but $b \notin \{d\}$.

    $\Rightarrow$ $\{d\}$ does \textbf{not} contain all the arguments it defends.

    $\Rightarrow$ $\{d\}$ is \textbf{not} complete.
\end{itemize}

\textbf{Conclusion:}
\[
\emptyset
\]
is the only complete extension.

\subsection{Exercise 5}

Attacks:
\[
b \to a,\quad c \to b,\quad d \to c
\]

\textbf{Grounded procedure (as in the slides):}

\begin{itemize}
    \item Step 0: all arguments are marked \textbf{undecided}

    \item Step 1: mark as \textbf{accepted (green)} any argument not attacked by a green or undecided argument

    The only such argument is:
    \[
    d
    \]

    \item Step 2: mark as \textbf{rejected (red)} any argument attacked by a green argument

    Since $d \to c$, we mark:
    \[
    c \text{ as rejected}
    \]

    \item Step 3: repeat Step 1

    Now, $b$ is attacked only by $c$, which is rejected, hence it is not attacked by any green or undecided argument.

    Therefore:
    \[
    b \text{ is accepted}
    \]

    \item Step 4: apply Step 2 again

    Since $b \to a$, we mark:
    \[
    a \text{ as rejected}
    \]
\end{itemize}

\textbf{Final labelling:}
\[
\text{accepted: } \{d, b\}, \quad
\text{rejected: } \{c, a\}
\]

\textbf{Grounded extension:}
\[
\{d, b\}
\]

\textbf{Conclusion:} the grounded extension is not empty.

\subsection{Exercise 6}

Attacks:
\[
a \leftrightarrow b,\quad c \to b
\]

\textbf{Complete extensions:}
\[
\emptyset,\ \{a\},\ \{a,c\}
\]

\textbf{Preferred (maximal complete):}
\[
\{a,c\}
\]

\textbf{Grounded:}
\[
\emptyset
\]

\textbf{Observation:}
Preferred $\neq$ grounded.

\subsection{Exercise 7}

Cycle:
\[
a \to b \to c \to a
\]

\textbf{Stable condition:}
Must attack all arguments outside the set.

Try candidates:
\begin{itemize}
    \item $\{a\}$: does not attack $c$
    \item $\{b\}$: does not attack $a$
    \item $\{c\}$: does not attack $b$
\end{itemize}

\textbf{Conclusion:}
No stable extension exists.

\subsection{Exercise 8}

Attacks:
\[
a \to b,\quad b \leftrightarrow c,\quad d \to a,\quad e \to c
\]

\textbf{Naive:}
\[
\{a,c\},\ \{d,b\},\ \{d,c\},\ \{a,e\}
\]

\textbf{Admissible (example):}
\[
\{d\},\ \{e\}
\]

\textbf{Complete (example):}
\[
\{d\}
\]

\textbf{Grounded:}
Start with unattacked:
\[
d, e
\]
Then propagate:
\[
\{d,e\}
\]

\textbf{Preferred (example):}
\[
\{d,e\}
\]

\textbf{Stable (check):}
\[
\{d,e\}
\]
attacks $a$ and $c$, but not $b$ $\Rightarrow$ not stable

Try $\{a,e\}$:
attacks $b$ and $c$ $\Rightarrow$ stable

\textbf{Conclusion:}
All semantics differ.

\end{document}